\section{Architecture (v1 + v2 split)}
\subsection{System split}
The repository is intentionally partitioned: locked v1 assets preserve the historical implementation baseline, while v2 hosts active experimental development. This separation supports auditability and controls accidental regression in legacy code.

\subsection{Core architectural elements}
\paragraph{Graph memory.}
Concept nodes and weighted relations form the persistent substrate. State snapshots are serialized to JSON for post hoc analysis.

\paragraph{Wave-state dynamics (ECWF).}
An ECWF layer models dynamic activation and coherence propagation, enabling non-static concept salience and context-sensitive interaction.

\paragraph{Governance (Three Kings).}
A governance layer regulates admission, stabilization, and policy-constrained transitions; this layer provides guardrails for concept promotion and retention.

\paragraph{Nine-stage pipeline.}
Verdant v2 organizes cognition-like updates in nine stages: ingest, parse, map, excite, cohere, evaluate, gate, consolidate, and persist. These stages provide deterministic checkpoints for instrumentation.

\subsection{Figure placeholders}
\begin{figure}[h]
  \centering
  \fbox{\parbox{0.9\linewidth}{\centering Placeholder: full architecture diagram (v1/v2 split and data flow).}}
  \caption{Verdant-Minds architecture and control surfaces.}
\end{figure}
